\chapter{Reductions, Concrete Bounds, and Paper Algebra}
\label{ch:bounds}

\section{Reduction Terms}

\code{HAETAE\_Reductions.ec} collects the algebra of the final security bound.
It defines and relates terms for signing-query scaling, ROM collision losses,
prequery losses, min-entropy losses, rejection-sampling losses, and bimodal to
Module-SIS reduction losses.

\section{Pen-and-Paper Bound Composition}

Let the game hops produce losses \(\epsilon_1,\ldots,\epsilon_k\), and let the
final extractor produce reductions \(B_1,\ldots,B_r\).  At the schematic
paper-theorem level, the standard game-hop composition gives
\[
\Adv^{\eufcma}_{\haetae}(A)
\le
\sum_{j=1}^r \Adv_j(B_j) + \sum_{i=1}^k \epsilon_i.
\]
For \haetae{}, the intermediate \(\epsilon_i\) ledger includes ROM collision
terms, programmed-site prequery terms, sampler-prequery terms, min-entropy
terms, rejection-sampling losses, and exact-sampling replacement losses.  This
is the paper-style hop ledger, not the final counted-ROM right-hand side.  In
the final checked implication, the rejection and exact-sampling obligations are
antecedents or bridge-lemma costs, while the displayed final bound adds only the
hardness terms, the bimodal-to-Module-SIS reduction term, and the counted ROM
collision/prequery/min-entropy terms.  This is a logical grouping statement, not
by itself a non-vacuity statement: with the current constants the
paper-style \ec{rejection\_sampling\_loss\_term} is larger than one, so any final
antecedent that adds this term to a probability must be separately checked for
satisfiability before the implication is cited as a useful concrete bound.
The checked instantiation discussed in this dissertation replaces the schematic
left-hand side by \(\Adv_{\mathsf{ECModel},\rom}^{\eufcma}(A)\), while the
right-hand-side bound ledger is instantiated at \(q_H=q_S=2\).

\begin{proof}[Pen-and-paper proof of composition]
Let \(G_0\) be the real \eufcma{} game and \(G_k\) be the final reduction game.
For each adjacent pair, prove either exact equality of the relevant event or a
statistical-distance bound
\(|\Pr[G_i]-\Pr[G_{i+1}]|\le\epsilon_i\).  By the triangle inequality,
\(|\Pr[G_0]-\Pr[G_k]|\le\sum_i\epsilon_i\).  The final game success is bounded
by the reduction advantages \(\sum_j\Adv_j(B_j)\).  Combining the inequalities
gives the result.
\end{proof}

\section{Nonnegativity and Cardinality}

Concrete bounds require side conditions.  The proof includes nonnegativity
lemmas for budgets and loss terms, lower bounds for challenge support, and
subunit bounds where probabilities must be at most one.  This arithmetic layer
is critical: it ensures the final implication composes valid inequalities rather
than relying on informal real-number reasoning.

\section{Detailed EasyCrypt Code Analysis}

\code{HAETAE\_Reductions.ec} should be read as the arithmetic ledger of the
security proof.  Lemmas ending in \code{E} are usually expansion lemmas: they
rewrite named loss terms into their concrete expressions.  Lemmas ending in
\code{nonnegative}, \code{positive}, \code{le1}, or \code{subunit} discharge
side conditions required by probability reasoning.

For example, \code{counted\_rom\_programming\_loss\_term\_concreteE} connects a
named loss to its concrete budget/support expression, while
\code{counted\_rom\_programming\_loss\_term\_subunit} proves that the expression
is a valid probability-scale quantity.  Similarly,
\code{challenge\_support\_cardinality\_lower\_bound\_positive} prevents a hidden
division-by-zero assumption in entropy-based bounds.

The theorem file \code{HAETAE\_Security.ec} consumes these arithmetic facts
rather than re-proving arithmetic at every hop.  This separation improves proof
maintenance: changing a bound expression should affect the reduction ledger and
a small number of bridge lemmas, not the entire game-hop development.

\section{Hardness Boundary}

The final proof is conditional on the assumptions encoded in
\code{HAETAE\_Assumptions.ec}.  The machine-checked result therefore states
that the modeled \haetae{} adversary's advantage is bounded by formal reduction
advantages and checked loss terms under those assumptions.
